\section{Discussion}\label{sec:discussion}

\subsection{What the demonstration establishes about the method}

The method was proposed as a way of making diagnosability a criterion in architecture selection. The demonstration supports four claims about it that do not depend on which machine came out ahead.

\emph{The design-level description of a failure mode does not fix its physical severity.} A short-circuit of the same winding fraction produced a fault current that differed by up to a factor of 1.3 between the PMSG and the SCIG and, relative to the current the stator sensors see, by a factor of 2.3 (Section~\ref{sec:res-severity}). In the WFSG benchmark the experimenters had raised the fault impedance with extent, so that on that bench the extent axis and the impedance axis are confounded. Any design method that compares alternatives on a failure-mode class must therefore carry two severity axes---the design-level one, on which requirements are written, and a physical one, on which the alternatives are compared---and must state which one a number refers to. The SDR statistic makes this explicit, because it can be plotted against either (Figs.~\ref{fig:severity} and \ref{fig:physical}).

\emph{Where a failure becomes observable is decided by the control structure, not by the fault.} With the fault physically similar, the PMSG's current controller suppressed the current asymmetry and pushed it into the terminal voltage, whereas in the SCIG the asymmetry appeared in the magnetising-axis current and its controller action, and in the WFSG additionally in the field current and the current harmonics. The two detectors trained on the PMSG and SCIG weight the same features almost orthogonally. This is the mechanism behind the headline design finding: the same monitoring specification (``three current transducers, current-signature analysis'') gave a minimum detectable inter-turn extent of 7.4\,\% on the PMSG and 2.8\,\% on the SCIG, and adding three voltage transducers closed the gap on the PMSG but bought nothing on the SCIG. The cheapest adequate observation suite is a property of the alternative, and the reason can be read off the signature-location statistic.

\emph{A monitoring design transfers across alternatives only under the right referencing rule.} A detector built on one machine and applied to another in physical units lost most of its performance; scaling its features on the target machine's healthy trials recovered part of it; referencing each feature to the trial's own healthy interval recovered most of it, in all six directions between the three machines and when two machines were pooled to detect faults in the third. For a product family this is a specification rule, not an algorithmic detail: the monitoring function should be defined on deviations from the unit's own reference behaviour, and the healthy trials that establish that reference are part of the commissioning procedure.

\emph{The measure is robust to its own tuning.} The null quantile and the window length moved the reported quantities in the expected direction and by small amounts (Section~\ref{sec:res-robust}); the bootstrap intervals separate the alternatives on the features where the signature relocates and overlap on the features where it does not. The method does not manufacture differences.

\subsection{Design implications}

For the wind-turbine generator, the results translate into statements a design team can use. A PMSG with full-power conversion should be specified with terminal voltage measurement if inter-turn faults below about 7\,\% of a winding are to be detected by threshold monitoring; controller-internal quantities are not a substitute, because the commanded voltage carries a weaker and less reliable version of the signature than the measured one. An SCIG reaches the same detectability with current transducers alone or with no added sensor at all, and additional hardware buys robustness rather than reach. A wound-field machine offers a further drive-internal channel, the field current, whose 2$f_e$ ripple detected every inter-winding fault and most inter-turn faults in the trials. A shaft torque transducer is not worth specifying for winding faults on any of the three. And a monitoring function meant to serve more than one generator type should be specified on self-referenced deviations.

More generally, the method gives DFMEA's detection rating a measured basis: the ``D'' of a failure mode becomes a function of the alternative and the observation suite, with an uncertainty interval, rather than a consensus value. It also gives digital-twin efforts a concrete design-stage use: a twin that can inject failure modes of controlled extent and produce the same observables as the monitoring suite is exactly the trial generator of step~2, and its fidelity can be checked with the same statistic against a small number of physical trials, as the three benches were used here.

\subsection{Applying the method elsewhere}

Nothing in Section~\ref{sec:method} is specific to electrical machines. What changes between domains is the content of the six framing decisions: what the alternatives are, what design-level extent scale is realisable in all of them, which suites are plausible, and what constitutes a trial. A gearbox family would use tooth-damage size as extent, vibration and oil-debris suites, and run-to-failure or seeded-fault trials; a battery-pack family would use cell-capacity loss, voltage/temperature/impedance suites, and cycling trials; a hydraulic actuator would use leakage cross-section, pressure/flow/position suites, and bench trials. The SDR, the pooled null, the reliability screen, the fixed/oracle distinction and the transfer test carry over unchanged. The one requirement the method imposes on the trial design---that every trial contain a healthy reference interval of at least twice the faulted duration---is mild and, as the WFSG benchmark shows, is met even by data collected for other purposes.

\subsection{Threats to validity}

The most important limitation is that each alternative is represented by one physical machine. The differences reported are differences between three machines; attributing them to topology rather than to the individual units requires either replicated machines or, more realistically at the design stage, simulated populations of each alternative under parameter variation, validated against the measured units. We regard the demonstration as establishing the method and the mechanisms, not the population-level ranking of the three topologies. A related confound is the sensing chain: the SCIG's magnetising current makes its stator current 2.5 times larger than the PMSG's on the same current transducers, and its relative healthy drift is correspondingly smaller (Table~\ref{tab:null}); part of the SCIG's advantage on current-side features may therefore be an effect of signal level on identical sensors rather than of topology. The method exposes this through the null table but cannot remove it without additional trials.

The benches are laboratory scale (2--2.5\,kVA); the operating points were scaled from a 15\,MW turbine, but the fault physics were not, and the fault durations were short (0.12--0.40\,s). The fault impedance was fixed for the PMSG and SCIG and confounded with extent for the WFSG. Only winding faults were studied, and only electrical and controller observables; the method is not restricted to either, but its behaviour on vibration or thermal observables and on mechanical faults remains to be shown. The WFSG has no healthy trials, so its reliability screen rests on the pre-fault null alone, and its operating points are unbalanced. The transfer experiment's negatives are dominated by reference and recovery windows; the fault-time windows of healthy trials, which contain the contactor operation without a fault, are the only negatives that resemble a nuisance transient, and other transients (load steps, grid events) were not in the data. Finally, the minimum detectable extent is a discrete quantity on a coarse extent grid; its bootstrap intervals are correspondingly wide, and the continuous SDR is the better quantity for comparing alternatives.
